\documentclass[12pt,a4paper]{article}
\usepackage[utf8]{vietnam}
\usepackage[top=2cm, bottom=2cm, left=1.5cm, right=1cm]{geometry}
\usepackage{amsmath,amsfonts,amssymb}
\usepackage{tabularx,longtable,multirow,makecell} % Thêm gói makecell
\usepackage{array}
\usepackage{fancyhdr}
\usepackage{enumitem}

% Cấu hình cột bảng
\newcolumntype{L}[1]{>{\raggedright\arraybackslash}p{#1}}
\newcolumntype{C}[1]{>{\centering\arraybackslash}p{#1}}

% Tăng độ giãn dòng để các ô không bị dính chữ
\renewcommand{\arraystretch}{1.6} 

% Cấu hình Header/Footer
\pagestyle{fancy}
\fancyhf{}
\rhead{\small Tổ Toán}
\lhead{\small Phân phối chương trình Toán 12 - Buổi 2}
\cfoot{\thepage}

\begin{document}

% --- PHẦN TIÊU NGỮ ---
\noindent
\begin{minipage}[t]{0.5\textwidth}
    \centering
    \textbf{TRƯỜNG THPT NGUYỄN HỮU CẢNH}\\
    \textbf{TỔ TOÁN}
\end{minipage}
\begin{minipage}[t]{0.5\textwidth}
    \centering
    \textbf{CỘNG HÒA XÃ HỘI CHỦ NGHĨA VIỆT NAM}\\
    \textbf{Độc lập – Tự do – Hạnh phúc}\\
    \rule{3cm}{0.4pt}
\end{minipage}

\vspace{0.8cm}

\begin{center}
    \Large \textbf{PHÂN PHỐI CHƯƠNG TRÌNH MÔN TOÁN – KHỐI 12} \\
    \large \textbf{Năm học 2025 – 2026} \\
    \large \textit{Buổi 2 Học kì II: 17 tuần}
\end{center}

\vspace{0.5cm}

% --- BẢNG PHÂN PHỐI CHI TIẾT ---
% Điều chỉnh lại độ rộng cột Tuần một chút để ngày tháng nằm gọn trên 2 dòng
\begin{longtable}{|C{1cm}|C{2.8cm}|C{2cm}|C{1.5cm}|L{8.7cm}|}
\hline
\textbf{STT} & \textbf{Tuần} & \textbf{Tên môn} & \textbf{TIẾT} & \textbf{Tên bài dạy và Yêu cầu cần đạt} \\ \hline
\endhead

% TUẦN 1
1 & \multirow{2}{*}{\makecell{1 \\ (19/01--25/01)}} & \multirow{2}{*}{Toán 12} & 1 B2 & \textbf{Chuyên đề nâng cao Nguyên hàm.} \\ \cline{1-1} \cline{4-5}
2 & & & 2 B2 & \textit{Yêu cầu:} Vận dụng các phương pháp đổi biến, từng phần tìm nguyên hàm hàm hợp. \\ \hline

% TUẦN 2
3 & \multirow{2}{*}{\makecell{2 \\ (26/01--01/02)}} & \multirow{2}{*}{Toán 12} & 3 B2 & \textbf{Chuyên đề nâng cao PT mặt phẳng.} \\ \cline{1-1} \cline{4-5}
4 & & & 4 B2 & \textit{Yêu cầu:} Giải bài toán cực trị về khoảng cách và góc liên quan đến mặt phẳng. \\ \hline

% TUẦN 3
5 & \multirow{2}{*}{\makecell{3 \\ (02/02--08/02)}} & \multirow{2}{*}{Toán 12} & 5 B2 & \textbf{Chuyên đề nâng cao tích phân.} \\ \cline{1-1} \cline{4-5}
6 & & & 6 B2 & \textit{Yêu cầu:} Tính tích phân hàm ẩn và ứng dụng giải bài toán thực tế. \\ \hline

\multicolumn{5}{|c|}{\textbf{NGHỈ TẾT NGUYÊN ĐÁN (09/02 $\to$ 22/02)}} \\ \hline

% TUẦN 4
7 & \multirow{2}{*}{\makecell{4 \\ (23/02--01/03)}} & \multirow{2}{*}{Toán 12} & 7 B2 & \textbf{Chuyên đề nâng cao PT mặt phẳng (tt).} \\ \cline{1-1} \cline{4-5}
8 & & & 8 B2 & \textit{Yêu cầu:} Vận dụng phương trình mặt phẳng giải các bài toán tương giao phức tạp. \\ \hline

% TUẦN 5
9 & \multirow{2}{*}{\makecell{5 \\ (02/03--08/03)}} & \multirow{2}{*}{Toán 12} & 9 B2 & \textbf{CĐNC Ứng dụng hình học của tích phân.} \\ \cline{1-1} \cline{4-5}
10 & & & 10 B2 & \textit{Yêu cầu:} Tính thiết diện, thể tích vật thể có hình dạng đặc biệt bằng tích phân. \\ \hline

% TUẦN 13 (Ví dụ tuần có 3 tiết để thầy thấy sự linh hoạt)
25 & \multirow{3}{*}{\makecell{13 \\ (27/04--03/05)}} & \multirow{3}{*}{Toán 12} & 25 B2 & \textbf{Ôn tập tốt nghiệp THPT.} \\ \cline{1-1} \cline{4-5}
26 & & & 26 B2 & \multirow{2}{=}{\textit{Yêu cầu:} Luyện giải các dạng toán thực tế và vận dụng cao trong đề thi chính thức.} \\ \cline{1-1} \cline{4-4}
27 & & & 27 B2 & \\ \hline

% CÁC TUẦN CUỐI (Sửa lỗi dính chữ ở ảnh thầy gửi)
28 & \makecell{14 \\ (04/05--10/05)} & Toán 12 & 28 B2 & \textbf{Ôn tập tốt nghiệp THPT.} \\ \hline

29 & \multirow{2}{*}{\makecell{15 \\ (11/05--17/05)}} & \multirow{2}{*}{Toán 12} & 29 B2 & \multirow{2}{=}{\textbf{Ôn tập cuối năm, ôn thi tốt nghiệp.}} \\ \cline{1-1} \cline{4-4}
30 & & & 30 B2 & \\ \hline

31 & \multirow{2}{*}{\makecell{16 \\ (18/05--24/05)}} & \multirow{2}{*}{Toán 12} & 31 B2 & \multirow{2}{=}{\textbf{Ôn tập cuối năm, ôn thi tốt nghiệp.}} \\ \cline{1-1} \cline{4-4}
32 & & & 32 B2 & \\ \hline

33 & \multirow{2}{*}{\makecell{17 \\ (25/05--30/05)}} & \multirow{2}{*}{Toán 12} & 33 B2 & \multirow{2}{=}{\textbf{Ôn tập cuối năm, ôn thi tốt nghiệp.}} \\ \cline{1-1} \cline{4-4}
34 & & & 34 B2 & \\ \hline

\end{longtable}

\vspace{1.0cm}
% --- CHỮ KÝ ---
\noindent
\begin{minipage}[t]{0.45\textwidth}
    \centering
    \textbf{P. HIỆU TRƯỞNG}\\
    \vspace{2.5cm}
    \textbf{...........................................}
\end{minipage}
\hfill
\begin{minipage}[t]{0.45\textwidth}
    \centering
    \textit{Tp. HCM, ngày 05 tháng 09 năm 2025}\\
    \textbf{TỔ TRƯỞNG}\\
    \vspace{2cm}
    \textbf{ĐỖ THỊ BẠCH LAN}
\end{minipage}

\end{document}